% Part: methods
% Chapter: induction
% Section: relations

\documentclass[../../../include/open-logic-section]{subfiles}

\begin{document}

\olfileid{mth}{ind}{rel}

\olsection{اړيکې او تابعې}

کله چې د څيزونو يو سټ، لکه طبيعي عددونه يا ښه‌جوړ
ترمونه، په استقرايي ډول تعريف کړو، نو \emph{پر دغو
څيزونو اړيکې} هم په استقرا تعريفولاے شو۔ د بېلګې
په توګه، دا نظر په پام کښې ونيسئ: يو ښه‌جوړ
ترم~$t_1$ د ښه‌جوړ ترم~$t_2$ فرعي ترم دے، که د
هغه د يوې برخې په توګه پکښې راځي۔ يوه نښه ورته
کاروو: $t_1 \sqsubseteq t_2$۔ البته هر ښه‌جوړ
ترم د خپل ځان فرعي ترم دے: $t \sqsubseteq t$۔
د دې اړيکې استقرايي تعريف داسې ورکولاے شو:

\begin{defn}
د ښه‌جوړ ترم~$t_1$ فرعي ترم کېدل د~$t_2$، يعنې
$t_1
  \sqsubseteq t_2$ اړيکه، پر~$t_2$ په استقرا
داسې تعريفېږي:
  \begin{enumerate}
  \item که $t_2$ يو تورے وي، نو $t_1 \sqsubseteq t_2$
    هله او يوازې هله چې $t_1 = t_2$۔
  \item که $t_2$، $[s_1 \circ s_2]$ وي، نو
    $t_1 \sqsubseteq t_2$ هله او يوازې هله چې
    $t_1 =
    t_2$، $t_1 \sqsubseteq s_1$، يا $t_1 \sqsubseteq s_2$۔
  \end{enumerate}
\end{defn}

د بېلګې په توګه، دا تعريف به راته ووايي چې
$\mathrm{a}
\sqsubseteq [\mathrm{b} \circ \mathrm{a}]$۔ ځکه چې
(2) وايي $\mathrm{a} \sqsubseteq [\mathrm{b} \circ \mathrm{a}]$
هله او يوازې هله چې $\mathrm{a}
= [\mathrm{b} \circ \mathrm{a}]$، يا $\mathrm{a} \sqsubseteq b$،
يا $\mathrm{a} \sqsubseteq \mathrm{a}$۔ لومړني دوه
ناسم دي: $\mathrm{a}$ ښکاره له $[\mathrm{b} \circ
  \mathrm{a}]$ سره يو شان نۀ دے، او د~(1) له مخې
$\mathrm{a} \sqsubseteq \mathrm{b}$ هله او يوازې
هله چې $\mathrm{a} = \mathrm{b}$، چې دا هم ناسم
دے۔ خو بيا هم د~(1) له مخې $\mathrm{a} \sqsubseteq \mathrm{a}$
هله او يوازې هله چې $\mathrm{a} = \mathrm{a}$،
چې دا رښتيا دے۔

دا مهمه ده چې پوه شو د دې تعريف برياليتوب پر
يو داسې حقيقت تکيه کوي چې تر اوسه مو نۀ دے
ثابت کړے: هر ښه‌جوړ ترم~$t$ يا يوازې يو تورے
دے، يا \emph{په يکتا ډول ټاکل شوي} ښه‌جوړ
ترمونه $s_1$ او $s_2$ شته چې $t = [s_1 \circ s_2]$۔
دلته «په يکتا ډول ټاکل شوي» دا معنا لري چې
که $t = [s_1 \circ s_2]$ وي، نو دا \emph{هم}
$= [r_1 \circ r_2]$ نۀ وي، له $s_1 \neq r_1$
يا $s_2 \neq r_2$ سره۔ که داسې کېداے، نو (2)
شرط ښايي له خپل ځان سره ټکر وکړي: $t_2$ چې
د $[s_1 \circ s_2]$ په توګه ولولو، ښايي
$t_1
\sqsubseteq t_2$ ترلاسه کړو، خو که $t_2$ د
$[r_1 \circ r_2]$ په توګه ولولو، ښايي
$t_1 \sqsubseteq t_2$ نۀ وي، دا نتيجه ترلاسه کړو۔ مخکې
له دې چې ثابت کړو دا نۀ شي کېداے، يوه داسې
بېلګه ګورو چې دا پکښې \emph{کېداے} شي۔

\begin{defn}
\emph{بې‌قوسه ترمونه} په استقرا داسې تعريف کړئ:
  \begin{enumerate}
  \item هر تورے يو بې‌قوسه ترم دے۔
    \item که $s_1$ او $s_2$ بې‌قوسه ترمونه وي، نو
      $s_1 \circ s_2$ بې‌قوسه ترم دے۔
    \item نور هېڅ څه بې‌قوسه ترم نۀ دي۔
  \end{enumerate}
\end{defn}

بې‌قوسه ترمونه، د بېلګې په توګه، $\mathrm{a}$،
$\mathrm{b} \circ
\mathrm{d}$، $\mathrm{b} \circ \mathrm{a} \circ \mathrm{b}$
دي۔ اوس که د بې‌قوسه ترمونو دپاره «فرعي ترم»
په هماغه ډول تعريف کړو لکه پورته مو چې وکړ،
دويم شرط به داسې وي:
\begin{center}
که $t_2 = s_1 \circ s_2$ وي، نو $t_1 \sqsubseteq t_2$
هله او يوازې هله چې $t_1 = t_2$، $t_1
  \sqsubseteq s_1$، يا $t_1 \sqsubseteq s_2$۔
\end{center}

اوس $\mathrm{b} \circ \mathrm{a} \circ \mathrm{b}$
د $s_1
\circ s_2$ په بڼه دے، په دې ډول چې:
\begin{align*}
  s_1 & = \mathrm{b} \text{ او} & s_2 & = \mathrm{a} \circ \mathrm{b}.
\intertext{دا د $r_1 \circ r_2$ په بڼه هم دے، په دې ډول چې:}
  r_1 & = \mathrm{b} \circ \mathrm{a} \text{ او} & r_2 &= \mathrm{b}.
\end{align*}
اوس آيا $\mathrm{a} \circ \mathrm{b}$ فرعي ترم دے
د $\mathrm{b} \circ
\mathrm{a} \circ \mathrm{b}$؟ که لومړے لوست
تعقيب کړو، ځواب هو دے، او که دويم تعقيب
کړو، ځواب نۀ دے۔

دا خاصيت چې له نورو ښه‌جوړو ترمونو څخه د يوه
ښه‌جوړ ترم د جوړېدو طريقه يکتا ده،
\emph{يکتا لوستونتيا} نومېږي۔ ځکه چې د داسې
په استقرايي ډول تعريف شوو څيزونو دپاره د
اړيکو استقرايي تعريفونه مهم دي، نو بايد
ثابت کړو چې دا خاصيت سم دے۔

\begin{prop}
فرض کړئ چې $t$ يو ښه‌جوړ ترم دے۔ بيا يا $t$
يوازې يو تورے دے، يا په يکتا ډول ټاکل شوي
ښه‌جوړ ترمونه $s_1$، $s_2$ شته چې~$t =
  [s_1 \circ s_2]$۔
\end{prop}

\begin{proof}
که $t$ يوازې يو تورے وي، شرط پوره دے۔ نو
فرض کړئ چې $t$ يوازې يو تورے نۀ دے۔ له
استقرايي تعريفه معلومولاے شو چې بيا $t$
خامخا د $[s_1 \circ s_2]$ په بڼه وي، د
ځينو ښه‌جوړو ترمونو $s_1$ او~$s_2$ دپاره۔
دا ثابتول پاتې دي چې په يکتا ډول ټاکل
شوي دي؛ يعنې که $t = [r_1 \circ r_2]$، نو
$s_1 = r_1$ او $s_2 = r_2$۔

نو فرض کړئ چې $t = [s_1 \circ s_2]$ او
هم $t = [r_1 \circ r_2]$، د ښه‌جوړو
ترمونو $s_1$، $s_2$، $r_1$، $r_2$
دپاره۔ بايد وښيو چې $s_1 = r_1$ او
$s_2 = r_2$۔ لومړے، $s_1$ او $r_1$
خامخا يو شان دي، ځکه که نۀ وي، نو
يو ئې د بل خاصه پيلنۍ برخه ده۔ خو
د \olref[sti]{prop:initial} له مخې دا
ناممکنه ده، که $s_1$ او~$r_1$ دواړه
ښه‌جوړ ترمونه وي۔ خو که $s_1 = r_1$
وي، نو ښکاره $s_2 = r_2$ هم دے۔
\end{proof}

تابعې هم په استقرايي ډول تعريفولاے شو۔
د بېلګې په توګه، تابع~$f$ چې هر ښه‌جوړ
ترم پکښې د يو بل دننه د راغلو $[\dots]$
قوسونو تر ټولو زيات ژوروالي ته لېږي،
داسې تعريفولاے شو:

\begin{defn}
\ollabel{defn:depth} د يوه ښه‌جوړ ترم
\emph{ژوروالے}، $f(t)$، په استقرايي
ډول داسې تعريفېږي:
  \[
  f(t) = \begin{cases}
    0 & \text{ که $t$ يو تورے وي}\\
    \max(f(s_1), f(s_2)) + 1 & \text{ که $t = [s_1 \circ s_2]$ وي۔}
  \end{cases}
  \]
\end{defn}

د بېلګې په توګه:
\begin{align*}
  f([\mathrm{a} \circ \mathrm{b}]) & =
    \max(f(\mathrm{a}),f(\mathrm{b})) + 1 = \\
  &= \max(0, 0) + 1 = 1, \text{ او}\\
  f([[\mathrm{a} \circ \mathrm{b}] \circ \mathrm{c}]) & =
    \max(f([\mathrm{a} \circ \mathrm{b}]), f(\mathrm{c})) + 1 = \\
  & = \max(1,0) + 1 = 2.
\end{align*}

البته دلته فرض کوو چې $s_1$ او $s_2$
ښه‌جوړ ترمونه دي، او له دې حقيقته
کار اخلو چې هر ښه‌جوړ ترم يا يو
تورے دے يا د $[s_1 \circ s_2]$ په
بڼه دے۔ بيا دا مهمه ده چې په دغه
بڼه يوازې په يوه طريقه کېداے شي۔
د دې د علت د ليدلو دپاره هغه
بې‌قوسه ترمونه بيا په نظر کښې
ونيسئ چې مخکې مو تعريف کړل۔
ورته «تعريف» به داسې وي:
\[
  g(t) =
  \begin{cases}
    0 & \text{ که $t$ يو تورے وي}\\
   \max(g(s_1), g(s_2)) + 1 & \text{ که $t = s_1 \circ s_2$ وي۔}
  \end{cases}
\]
اوس بې‌قوسه ترم $\mathrm{a} \circ \mathrm{b} \circ
\mathrm{c} \circ \mathrm{d}$ په نظر
کښې ونيسئ۔ دا له يوې څخه په زياتو
طريقو لوستلاے شو؛ د بېلګې په توګه
د $s_1 \circ s_2$ په توګه، په دې
ډول چې:
\begin{align*}
  s_1 & = \mathrm{a} \text{ او} &
  s_2 & = \mathrm{b} \circ \mathrm{c} \circ \mathrm{d},
\intertext{يا د $r_1 \circ r_2$ په توګه، په دې ډول چې:}
  r_1 & = \mathrm{a} \circ b \text{ او} &
  r_2 &= \mathrm{c} \circ \mathrm{d}.
\end{align*}
د لومړي لوست له مخې د $g$ حساب
به راکړي:
\begin{align*}
  g(s_1 \circ s_2) & =
  \max(g(\mathrm{a}), g(\mathrm{b} \circ \mathrm{c} \circ \mathrm{d})) + 1 =\\ & =
  \max(0,2) + 1 = 3
  \intertext{خو د بل لوست له مخې ترلاسه کوو:}
  g(r_1 \circ r_2) & =
  \max(g(\mathrm{a} \circ \mathrm{b}), g(\mathrm{c} \circ \mathrm{d})) + 1=\\ &
  = \max(1,1) + 1 = 2
\end{align*}
خو يوه تابع بايد تل يو يکتا ارزښت
ورکړي؛ نو د~$g$ دغه «تعريف» هېڅ
تابع نۀ تعريفوي۔

\begin{prob}
د تابع $l$ استقرايي تعريف ورکړئ،
چې $l(t)$ د نښو شمېر دے په
ښه‌جوړ ترم~$t$ کښې۔
\end{prob}

\begin{prob}
پر ښه‌جوړو ترمونو $t$ په جوړښتي
استقرا ثابت کړئ چې $f(t) < l(t)$،
چې $l(t)$ د نښو شمېر دے په~$t$
کښې او $f(t)$ د $t$ ژوروالے دے،
لکه چې په \olref[mth][ind][rel]{defn:depth}
کښې تعريف شوے دے۔
\end{prob}

\end{document}
